% ===================================================================
%  TABULKA č. 7 — Vesnice v zimě. — Selský dům na jaře.
%  Traduction 2026 d'apres texto/io, controlee sur texto/fr, texto/ru
%  et texto/uk. Consigne : voir texto/cs/10-tabulka-01.tex.
%
%  OUVERTURE DE LA DEUXIEME SERIE : « DRUHÁ SÉRIE ».
%
%  LE TABLEAU DE LA FERME EST CELUI OU LE TCHEQUE EST LE PLUS
%  LUI-MEME, comme il l'a ete pour les six colonnes precedentes et
%  pour la meme raison : la ferme s'apprend dehors et sans maitre.
%  « Radlice » le soc, « brázda » le sillon, « hrábě » le rateau,
%  « vidle » la fourche, « chomout » le collier, « chlév » l'etable,
%  « holubník » le colombier, « kurník » le poulailler, « hnůj » le
%  fumier, « stodola » la grange, « máselnice » la baratte.
%
%  ET LA MESURE DEGAGEE EN GALICIEN SE VERIFIE UNE SECONDE FOIS SUR
%  UN AUTRE COUPLE. La distance du tcheque a l'ukrainien se creuse a
%  mesure qu'on s'eloigne de l'ecole : elle est maximale ici et au
%  corps du tableau 2, minimale a la maladie du tableau 4. Deux
%  couples de voisines, deux familles differentes, une seule courbe.
% ===================================================================

%%K t07-noto-1 noto
\VUnotes{143.2mm}{%
(*) Podrobnosti o jídle viz tabulky 6 a 13.%
}

%%K t07-apar-1 apar
\VUsaut{4.0mm}
\VUcentre{13.2pt}{90}{DRUHÁ SÉRIE}
\VUsaut{3.0mm}
\VUfilet{16mm}
\VUsaut{5.6mm}
\VUtitre{15.0pt}{60}{TABULKA č. 7}
\VUsaut{3.0mm}

%%K t07-tit-1 sub
\VUcentre{12.0pt}{0}{\textit{První scéna.}}
\VUsaut{3.0mm}

%%K t07-tit-2 sub
\VUpk{10.2pt}{40}{Vesnice v zimě.}
\VUsaut{4.0mm}

%%K t07-c1-01-1 p
1. --- O novoročních prázdninách jsem doprovázel svého otce do Nového Města, do
vesničky, kde měl vyřídit několik obchodních záležitostí.

%%K t07-c1-02-1 p
2. --- Použili jsme \VUgras{dostavník} \textsuperscript{(11)}, který jezdí mezi městem
a tím městečkem; ale protože byla velká zima, nebyla ta cesta v málo pohodlném voze
příliš příjemná.

%%K t07-c1-03-1 p
3. --- Proto jsme pocítili opravdovou radost, když jsme viděli \VUgras{kočího}
zastavit svůj vůz na náměstí, před \VUgras{hostincem} \textsuperscript{(2)} \textit{U
Bílého medvěda}. \VUgras{Hostinský} \textsuperscript{(4)} nás očekával na prahu svého
domu a klidně kouřil svou \VUgras{dýmku.}

%%K t07-c1-03-2 p
Pozval nás dovnitř a já jsem se nedal pobízet, abych se usadil před ohněm z roští,
který praskal v ohništi. Tehdy jsem se velmi vysmíval ostrému \VUgras{severnímu
větru}, který skřípal starým \VUgras{vývěsním štítem} \textsuperscript{(3)} a odnášel
v černých vírech \VUgras{kouř} \textsuperscript{(1)} vycházející z komína.

%%K t07-c1-04-1 p
4. --- Byla hodina oběda. Bez dalšího čekání jsme se dobře najedli prostého, ale
chutného jídla, které nám podali (*).

%%K t07-tit-3 sub
\VUpk{10.2pt}{40}{Několik návštěv.}
\VUsaut{3.0mm}

%%K t07-c1-05-1 p
5. --- Po obědě jsem doprovázel svého otce, který je pojišťovacím jednatelem, k jeho
klientům. Nejprve jsme navštívili \VUgras{kováře} \textsuperscript{(5)}, který bydlí
blízko hostince. Byl zaměstnán kováním koně. Obdivoval jsem sílu jeho druha, který
pevně držel na své kožené zástěře \VUgras{kopyto} koně, zatímco kovář připevňoval
\VUgras{podkovu} \textsuperscript{(6)} silnými údery kladiva.

%%K t07-c1-06-1 p
6. --- Potom jsme šli k vesnickému \VUgras{ševci} \textsuperscript{(7)}, ke starému
Jakubovi. Ačkoli má za vývěsní štít překrásnou \VUgras{botu,} spokojoval se s tím, že
podrážel starý pár prošlapaných bot.

%%K t07-c1-06-2 p
Stařec nám řekl se smíchem: «Ve městě jsem byl „výrobce bot“ a neměl jsem zákazníky.
Zde jsem „opravář bot“ a práce je mnoho.»

%%K t07-c1-07-1 p
7. --- Mluvil nám o svém sousedu \VUgras{holiči} \textsuperscript{(8)}, s nímž
zákazníci nejsou spokojeni. Jeho \VUgras{břitvy} jsou stále vyštíplé a jeho
\VUgras{nůžky} nikdy dobře nestříhají.

Ale protože je jediným holičem ve vesnici, je člověk ovšem nucen dát se od něho holit
a stříhat. Ostatně je to lakomý a nevlídný muž.

%%K t07-c1-08-1 p
8. --- Naštěstí se mu ostatní \VUgras{sousedé} nepodobají. Například \VUgras{kolář}
\textsuperscript{(9)} je dobrotivý muž, pracovitý a ochotný. Je potěšením dát si od něho
opravit vůz. Jeho práce je vždy dobře dokončena.

%%K t07-c1-09-1 p
9. --- Starý Jakub si také nestěžoval na \VUgras{sedláře} \textsuperscript{(10)},
kterému někdy pomáhá šít \VUgras{postroje,} když práce spěchá.

%%K t07-c1-09-2 p
Můj otec se ho vyptával na jeho rodinu: «Všichni se mají dobře. Moje vnučka je ve
\VUgras{škole} \textsuperscript{(12)}. Její \VUgras{učitelka} \textsuperscript{(13)} je
s ní velmi spokojena, je dobrá \VUgras{žákyně} \textsuperscript{(14)}. Nepodobá se mému
darebnému synovi; ten myslí jen na hraní. \VUgras{Učitel} \textsuperscript{(15)} ho
nedokáže nic naučit.»

%%K t07-tit-4 sub
\VUsaut{3.0mm}
\VUpk{10.2pt}{40}{Vesnické klevety.}
\VUsaut{3.0mm}

%%K t07-c1-10-1 p
10. --- Stařec nám potom vypravoval novinky z vesnice. Bude se stavět nová škola,
neboť ta, která je zřízena v \VUgras{obecním domě}, se stala příliš malou. Ostatně to
je snadné, neboť postačí koupit část zahrady \VUgras{mlynáře.} To bude pro chudáka
velmi výhodné. Kvůli mrazu už \VUgras{mlýnské kameny} \VUgras{mlýna}
\textsuperscript{(16)} nemohou mlít obilí, protože \VUgras{kolo} \textsuperscript{(17)}
bylo zastaveno \VUgras{ledem} \textsuperscript{(25)} \VUgras{potoka}
\textsuperscript{(23)}.

%%K t07-c1-11-1 p
11. --- «Když mluvíme o stavbách, viděl jste náš nový \VUgras{kostel}
\textsuperscript{(18)}? Je velmi malebný pod svým sněhovým pláštěm. \VUgras{Vrány}
\textsuperscript{(19)}, které už nepoznávají svou starou zvonici, smutně sedají na
velkém stromě na \VUgras{hřbitově} \textsuperscript{(20)}.»

%%K t07-c1-12-1 p
12. --- Ta myšlenka na hřbitov připomněla ševci osoby, které můj otec znal a které
loni zemřely. «Kolik \VUgras{hrobů} \textsuperscript{(21)}, můj dobrý pane, přibylo k
ostatním pod \VUgras{cypřišem} \textsuperscript{(22)}! Smrt pokosila našeho starého
notáře. Všichni vesničané byli na jeho \VUgras{pohřbu.}»

%%K t07-c1-13-1 p
13. --- «Zde je jeho nástupce, který bruslí na potoce. Právě mi přišel dát opravit
řemínek své \VUgras{brusle} \textsuperscript{(26)}, který byl přetržený.»

%%K t07-c1-14-1 p
14. --- «Zde je také, na břehu potoka, nový \VUgras{polní hlídač}
\textsuperscript{(27)}. Je to bývalý četník, který postrachem naplňuje vesnické
darebáky. Utíkají hned, jakmile spatří jeho \VUgras{kamaše} \textsuperscript{(29)} a
jeho \VUgras{čepici} \textsuperscript{(28)}.»

%%K t07-c1-15-1 p
15. --- «Ha! Zde je starý Josef, který veze \VUgras{vozík otepí dříví}
\textsuperscript{(30)} pro pekaře. --- To je pravda --- řekl můj otec ---; poznávám
jeho spřežení \VUgras{mezků} \textsuperscript{(31)}. Potřebuji s ním mluvit, využiji
příležitosti. Na shledanou, otče Jakube.»

%%K t07-c1-16-1 p
16. --- Právě když jsme vycházeli, vesnické děti opustily školu a rozběhly se na
\VUgras{klouzačku} \textsuperscript{(33)} s nebezpečím, že upadnou a zlomí si při
\VUgras{pádu} \textsuperscript{(34)} nějaký úd. Jeden z nich si vzal své \VUgras{sáně}
\textsuperscript{(32)} u koláře, posadil na ně jednoho ze svých kamarádů a rozběhl se.

%%K t07-c1-17-1 p
17. --- Jiní se vrhli na \VUgras{sněhovou hromadu} \textsuperscript{(37)} blízko
\VUgras{mostu} \textsuperscript{(40)} a začali bombardovat \VUgras{sněhuláka}
\textsuperscript{(39)}, kterého postavili včera a kterému nasadili klobouk, dýmku a
brašnu přes rameno.

%%K t07-c1-18-1 p
18. --- Málo dbají na mrazivý vítr a na zimu. Většina nemá ani \VUgras{šálu}
\textsuperscript{(35)} a, aby se po boji \VUgras{sněhovými koulemi}
\textsuperscript{(38)} zase zahřáli, stačí jim hrst velmi horkých \VUgras{kaštanů}
\textsuperscript{(42)} koupených od staré \VUgras{prodavačky} Jakobíny, která se třese
zimou pod svým příliš tenkým \VUgras{šátkem} \textsuperscript{(43)}.

%%K t07-tit-5 sub
\VUsaut{3.0mm}
\VUcentre{12.0pt}{0}{\textit{Druhá scéna.}}
\VUsaut{3.0mm}

%%K t07-tit-6 sub
\VUpk{10.2pt}{40}{Selský dům na jaře.}
\VUsaut{4.0mm}

%%K t07-c2-01-1 p
1. --- Přes všechny radosti zimy zdravíme s hlubokou radostí návrat \VUgras{jara.}

%%K t07-c2-01-2 p
Jak veselý obraz představuje selský dvůr večer v měsíci květnu!

%%K t07-c2-02-1 p
2. --- Na obzoru \VUgras{slunce} \textsuperscript{(1)} pomalu zapadá a zaplavuje
\VUgras{venkov} \textsuperscript{(3)} svými purpurovými \VUgras{paprsky}
\textsuperscript{(2)}. Pilný \VUgras{oráč} \textsuperscript{(6)} spěchá zorat své
\VUgras{pole} \textsuperscript{(4)}.

%%K t07-c2-02-2 p
Svým \VUgras{pluhem} \textsuperscript{(5)} zapřaženým do statného \VUgras{koně}
\textsuperscript{(7)} kreslí \VUgras{brázdu} \textsuperscript{(8)}, do níž
\VUgras{rozsévač} \textsuperscript{(9)} plnou hrstí vhodí \VUgras{semeno,} které vydá
zlaté klasy.

%%K t07-c2-03-1 p
3. --- \VUgras{Sekáči} \textsuperscript{(11)} opatření svými kosami posekali trávu
louky, obraceči usušili \VUgras{seno} \textsuperscript{(10)} obracejíce je
\VUgras{vidlemi} \textsuperscript{(12)} a hráběmi, a hle, vracejí se.

%%K t07-c2-03-2 p
Jedni kráčejí před těžkým vozem taženým voly; nesou své nářadí na rameni. Jiní,
lenivější, se pohodlně usadili na vonném seně.

%%K t07-c2-04-1 p
4. --- Cestou přátelsky zdraví \VUgras{zahradníka} \textsuperscript{(16)}
zaměstnaného v \VUgras{ovocném sadu} \textsuperscript{(13)} prořezáváním
\VUgras{ovocných stromů} a roubováním \VUgras{hrušní} \textsuperscript{(14)}.
\VUgras{Švestky} \textsuperscript{(15)} začínají kvést a v dobře pohnojené
\VUgras{zelinářské zahradě} \textsuperscript{(17)} se zelenině, zelí a salátům už
dobře daří.

%%K t07-c2-04-2 p
Na zídce krásný \VUgras{páv} \textsuperscript{(18)} rozevírá svůj ocas, aby dal
obdivovat své překrásné peří.

%%K t07-tit-7 sub
\VUpk{10.2pt}{40}{Selský dvůr.}
\VUsaut{3.0mm}

%%K t07-c2-05-1 p
5. --- Dveře \VUgras{stáje} \textsuperscript{(19)} jsou otevřené a je vidět žlab a
\VUgras{podestýlku} \textsuperscript{(23)} \VUgras{klisny} \textsuperscript{(21)},
kterou \VUgras{podkoní} \textsuperscript{(20)} právě vypřáhl. \VUgras{Hříbě}
\textsuperscript{(22)}, tak komické se svýma dlouhýma nohama, poskakujíc následuje svou
matku.

%%K t07-c2-06-1 p
6. --- Blízko \VUgras{studny} \textsuperscript{(29)} jdou \VUgras{krávy}
\textsuperscript{(25)} pít do dřevěného \VUgras{koryta} \textsuperscript{(26)}, které
jim slouží za napajedlo. \VUgras{Tele} \textsuperscript{(24)} využívá příležitosti k
sání a \VUgras{býk} \textsuperscript{(27)}, čichaje dobrou vůni stohu, vesele bučí a
hrdě zvedá hlavu s mocnými \VUgras{rohy} \textsuperscript{(28)}.

%%K t07-c2-07-1 p
7. --- Všichni pracují. \VUgras{Služka} \textsuperscript{(32)} drží v ruce
\VUgras{lano} \textsuperscript{(31)} studny. Čerpá vodu \VUgras{vědrem}
\textsuperscript{(30)}, aby napojila dobytek. Blízko \VUgras{stodoly}
\textsuperscript{(33)} jakýsi muž hrabe slámu a přede dveřmi \VUgras{stavení}
\textsuperscript{(34)} stará \VUgras{přadlena} \textsuperscript{(35)} svým kolovratem
a svou \VUgras{přeslicí} přede \VUgras{len} nebo \VUgras{konopí} jako za dávných časů.

%%K t07-tit-8 sub
\VUsaut{3.0mm}
\VUpk{10.2pt}{40}{Sedlák.}
\VUsaut{3.0mm}

%%K t07-c2-08-1 p
8. --- \VUgras{Sedlák} \textsuperscript{(36)} řídí všechny ty práce a dává pokyny svým
čeledínům. Doporučuje uklidit pod střechu \VUgras{brány} \textsuperscript{(40)} a
\VUgras{motyku} \textsuperscript{(39)}, které zapomněli blízko \VUgras{boudy}
\textsuperscript{(38)} \VUgras{psa} \textsuperscript{(37)}.

%%K t07-c2-09-1 p
9. --- Jeho volání plaší \VUgras{holuba} \textsuperscript{(41)}, který letí ze všech
sil do \VUgras{holubníku} \textsuperscript{(42)}, a \VUgras{slepice}
\textsuperscript{(43)}, které spěchají zpátky do \VUgras{kurníku}
\textsuperscript{(44)}.

%%K t07-c2-10-1 p
10. --- Před \VUgras{ovčínem} \textsuperscript{(48)}, zde je starý \VUgras{pastýř}
\textsuperscript{(45)}, který přivádí zpět své \VUgras{stádo} \textsuperscript{(49)}.
Drže svou \VUgras{hůl} \textsuperscript{(46)}, která mu nikdy nechybí, stejně jako jeho
vybledlá \VUgras{halena} \textsuperscript{(47)}, vede do chléva \VUgras{skopce}
\textsuperscript{(50)}, \VUgras{ovce} \textsuperscript{(53)}, \VUgras{jehňata}
\textsuperscript{(54)}, \VUgras{kozy} \textsuperscript{(51)} a \VUgras{kůzlata}
\textsuperscript{(52)}. Jeho věrný \VUgras{pes} \textsuperscript{(55)} Argus přichází
poslední a svým štěkotem pobízí k pilnosti ty, kteří se opožďují.

%%K t07-c2-11-1 p
11. --- Všechen ten hluk a ten pohyb neruší klid dobré \VUgras{dojné krávy}
\textsuperscript{(56)}, která cinká svým \VUgras{zvonkem} \textsuperscript{(57)},
zatímco ji dojí přede dveřmi \VUgras{chléva} \textsuperscript{(58)}.

%%K t07-c2-12-1 p
12. Zdá se, že chápe, že musí dát své mléko, aby \VUgras{selka}
\textsuperscript{(60)} mohla připravit \VUgras{máslo} v \VUgras{máselnici}
\textsuperscript{(61)} nebo výborné \VUgras{sýry} \textsuperscript{(64)}, které
odkapávají z desky položené na \VUgras{kádi} \textsuperscript{(63)}.

%%K t07-tit-9 sub
\VUsaut{3.0mm}
\VUpk{10.2pt}{40}{Drůbežárna.}
\VUsaut{3.0mm}

%%K t07-c2-13-1 p
13. --- Celou \VUgras{mléčnici} \textsuperscript{(59)} udržuje ve velké čistotě
\VUgras{čeledín} \textsuperscript{(65)}, který právě umyl \VUgras{konvice na mléko}
\textsuperscript{(62)} ve velkém vědru, které nese v ruce.

%%K t07-c2-14-1 p
14. --- Se svými tlustými dřeváky kráčí poněkud těžkopádně, a tak plaší
\VUgras{krocana} \textsuperscript{(68)}, \VUgras{kachny} \textsuperscript{(66)},
\VUgras{kačery} \textsuperscript{(67)} a \VUgras{husy} \textsuperscript{(69)}, které
široce otvírají \VUgras{svůj zobák} \textsuperscript{(70)} a neklidně křičí.

%%K t07-c2-14-2 p
\VUgras{Kohout} \textsuperscript{(71)}, bojovnější, zvedá svůj červený \VUgras{hřebínek}
\textsuperscript{(72)}, potřásá peřím svého \VUgras{ocasu} \textsuperscript{(73)} a
nechává zaznít zvučné «kykyryký».

%%K t07-c2-15-1 p
15. --- \VUgras{Slepice matka} \textsuperscript{(74)} hrabe na \VUgras{hnoji}
\textsuperscript{(81)} se svými \VUgras{kuřátky} \textsuperscript{(75)}.
\VUgras{Sele} \textsuperscript{(78)} utíká ke své matce, obrovské \VUgras{svini}
\textsuperscript{(77)}, kterou vidíme blízko chlívku.

%%K t07-c2-16-1 p
16. --- Ve svých kotcích jedí \VUgras{králíci} \textsuperscript{(79)} lačně jemnou
\VUgras{trávu} \textsuperscript{(80)}, kterou jim přináší sedlákova dcera. Janička má
velmi ráda své králíky a každý den, když si došla pro čerstvá \VUgras{vejce}
\textsuperscript{(76)} do kurníku, přichází navštívit své přátelíčky.
